\chapter{ARSITEKTUR SISTEM}

\section{Diagram Arsitektur}
Gambar~\ref{fig:m2-architecture} menunjukkan tiga jalur utama. Jalur masuk membawa data dari sumber ke jurnal milik Aggregator. Jalur baca membawa permintaan client melalui Client API menuju Aggregator. Jalur distribusi meneruskan perubahan yang telah tersimpan ke dashboard dan penerima peringatan. Dengan susunan ini, tiap layanan mengakses data melalui layanan pemiliknya.

\begin{figure}[htbp]
\centering
\resizebox{\linewidth}{!}{%
\begin{tikzpicture}[
  box/.style={draw=archink,rounded corners=3pt,fill=archlight,align=center,
             inner sep=5pt,font=\fontsize{10}{12}\selectfont},
  side/.style={box,draw=archink!45,fill=white,align=left,text width=3.1cm},
  eventbox/.style={box,draw=archevent,text width=3.05cm},
  flow/.style={-{Latex[length=2mm]},draw=archink,semithick},
  event/.style={flow,draw=archevent},
  note/.style={font=\footnotesize,fill=white,inner sep=1.1pt,align=center}
]
  \node[box,text width=4.25cm] (users) at (6.9,1.2)
    {Media, Tim Lapangan,\\BNPB internal};
  \node[box,minimum width=2.8cm] (auth) at (1.9,1.2) {Auth\\Login dan refresh};
  \node[box,text width=4.25cm] (api) at (6.9,-0.55) {Client API dan Analitik};
  \node[side,align=center] (audit) at (12.0,-0.55) {Auditor\\Token Audit};
  \node[side] (sources) at (1.9,-2.7)
    {\textbf{Sumber data}\\BMKG dan PVMBG\\Sensor Pantai\\Feed Satelit (SSE)};
  \node[box,text width=4.25cm] (agg) at (6.9,-2.7)
    {\textbf{Aggregator} ($\times 3$)\\Jurnal, pemetaan,\\Audit API dan relay outbox};
  \node[eventbox] (broker) at (12.0,-2.7) {\textbf{JetStream} ($\times 3$)};
  \node[side] (ops) at (1.9,-5.35)
    {\textbf{Observabilitas}\\Prometheus, Grafana\\Alloy, Loki, Tempo};
  \node[box,text width=4.25cm] (db) at (6.9,-5.0)
    {PostgreSQL\\4 shard\\3 salinan per shard};
  \node[eventbox] (consumer) at (12.0,-4.55) {Dashboard Consumer\\Notifier Consumer};
  \node[eventbox] (receiver) at (12.0,-6.3) {Penerima Tim Lapangan};

  \draw[flow] (users.west) -- (auth.east);
  \draw[flow] (users) -- node[note,right]{HTTP + JWT} (api);
  \draw[flow] (api) -- node[note,right]{HTTP internal} (agg);
  \draw[flow] (audit.south west) -- ++(0,-0.36) -| ([xshift=1.75cm]agg.north);
  \draw[flow] (sources.east) -- (agg.west);
  \draw[flow] (agg) -- node[note,right]{SQL / TLS} (db);
  \draw[event] (agg.east) -- (broker.west);
  \draw[event] (broker) -- (consumer);
  \draw[event] (consumer) -- (receiver);
  \draw[flow,dashed,draw=archink] (agg.south west) -- (ops.east);
\end{tikzpicture}%
}
\caption{Jalur data pada M2. Panah berwarna mengikuti distribusi event. Dari kedua consumer, hanya Notifier yang meneruskan peringatan pesisir kritis ke Penerima Tim Lapangan. Garis putus-putus menunjukkan pengiriman telemetri.}
\label{fig:m2-architecture}
\end{figure}


\section{Komponen dan Tanggung Jawab}

Auth menerbitkan token, sedangkan Client API memverifikasinya secara lokal dan menerapkan hak akses pada setiap permintaan. Replika Aggregator ditemukan melalui satu alamat layanan. Untuk tiap sumber, hanya pemilik sesi yang mengambil data dan melanjutkan jurnal. Kunci kepemilikan dan nomor generasi berada di PostgreSQL agar replika yang mengambil alih dapat melanjutkan dari posisi yang tersimpan.

Aggregator menyimpan data tiap sumber dan perubahan kanonik pada unit penyimpanan yang telah ditentukan. Setiap perubahan kanonik yang perlu dipublikasikan dicatat di \textit{outbox} dalam transaksi yang sama. Relay membaca catatan tersebut dan menerbitkannya ke JetStream. Dashboard Consumer dan Notifier Consumer memakai posisi baca serta penyimpanan SQLite masing-masing. Notifier meneruskan peringatan pesisir kritis ke Penerima Tim Lapangan, yang menyimpan tanda terimanya sendiri. Dasbor Analitik pada Client API menghitung data melalui Aggregator dan berbeda fungsi dari Dashboard Consumer yang mengikuti arus event.

Audit API berada di Aggregator dengan token tersendiri. Auditor dapat membaca data, checksum, lokasi salinan, dan status masukan yang belum selesai. Jalur operasi BNPB untuk memeriksa dan mengulang pesan gagal tetap melalui Client API. Pemisahan hak akses ini menjaga agar kredensial auditor hanya berlaku untuk pemeriksaan.

\section{Batas Container, Penyimpanan, dan Protokol}
Gambar~\ref{fig:m2-containers} memperinci jalur aplikasi pada Gambar~\ref{fig:m2-architecture}. Aggregator mengakses data kanonik melalui PostgreSQL, sedangkan Client API dan auditor mengaksesnya lewat HTTP. Basis data internal Auth dan consumer terpasang sebagai volume masing-masing. Broker juga memakai volume sendiri untuk menyimpan stream dan posisi baca.

\begin{figure}[H]
\centering
\newcommand{\ownedvolume}[1]{%
  {\setlength{\fboxsep}{2pt}\fcolorbox{archink!60}{white}{\fontsize{8}{10}\selectfont #1}}%
}
\resizebox{\linewidth}{!}{%
\begin{tikzpicture}[
  component/.style={draw=archink,rounded corners=2pt,fill=archlight,align=center,
    minimum height=10mm,inner sep=4pt,font=\fontsize{10}{12}\selectfont},
  external/.style={component,dashed,fill=white},
  wire/.style={-{Straight Barb[length=2mm,width=1.6mm]},draw=archink,semithick},
  event/.style={wire,draw=archevent},
  label/.style={fill=white,inner sep=1pt,font=\fontsize{8}{10}\selectfont,align=center}
]
  \node[external,text width=4cm] (clients) at (7.15,0) {Media, Tim Lapangan,\\BNPB internal};

  \node[component,text width=3.65cm,minimum height=16mm,anchor=north] (auth) at (12.9,0.58)
    {\textbf{Auth}\\[1mm]\ownedvolume{SQLite \texttt{auth-data}}};

  \node[component,text width=4.0cm] (api) at (7.15,-1.9)
    {\textbf{Client API}\\Dasbor Analitik};
  \node[external,text width=3.65cm] (audit) at (12.9,-1.9)
    {Auditor\\HTTP + token Audit};

  \node[component,text width=3.35cm] (bmkg) at (2.0,0)
    {\textbf{BMKG mock}\\HTTP :8081};
  \node[component,text width=3.35cm] (pvmbg) at (2.0,-1.8)
    {\textbf{PVMBG mock}\\HTTP :8082};
  \node[component,text width=3.65cm,minimum height=16mm,anchor=north] (pantai) at (2.0,-2.9)
    {\textbf{Sensor Pantai}\\HTTP :8080\\[1mm]\ownedvolume{PVC \texttt{sensor-pantai-data}}};
  \node[external,text width=3.35cm] (feed) at (2.0,-5.7)
    {\textbf{Feed Satelit}\\SSE melalui HTTPS};

  \node[component,text width=4.0cm,minimum height=19mm] (agg) at (7.15,-4.2)
    {\textbf{Aggregator} ($\times$3)\\Ingest, Audit API, dan relay};
  \node[component,text width=3.65cm,minimum height=20mm,anchor=north] (database) at (12.9,-3.25)
    {\textbf{PostgreSQL}\\4 shard $\times$ 3 salinan\\[1mm]\ownedvolume{Data milik Aggregator}};

  \node[component,text width=4.0cm,minimum height=16mm,anchor=north] (nats) at (7.15,-6.02)
    {\textbf{NATS JetStream} ($\times$3)\\[1mm]\ownedvolume{PVC \textit{file store} JetStream}};

  \node[component,text width=3.35cm,minimum height=16mm,anchor=north] (dash) at (2.0,-8.52)
    {\textbf{Dashboard Consumer}\\[1mm]\ownedvolume{SQLite \texttt{dashboard-data}}};
  \node[component,text width=4.0cm,minimum height=16mm,anchor=north] (notify) at (7.15,-8.52)
    {\textbf{Notifier Consumer}\\[1mm]\ownedvolume{SQLite \texttt{notifier-data}}};
  \node[component,text width=3.65cm,minimum height=16mm,anchor=north] (receiver) at (12.9,-8.52)
    {\textbf{Penerima Tim Lapangan}\\[1mm]\ownedvolume{SQLite \texttt{field-alert-data}}};

  \draw[wire] (clients.east) -- node[label]{HTTP} (auth.west |- clients.east);
  \draw[wire] (clients) -- node[label,right]{HTTP + JWT} (api);
  \draw[wire] (api) -- node[label,left]{HTTP internal} (agg);
  \draw[wire] (audit.west) -- (10.12,-1.9) |- ([yshift=5.5mm]agg.east);
  \draw[wire] (agg.east) -- node[label,above]{SQL/TLS} (database.west |- agg.east);

  \draw[wire] ([yshift=6.6mm]agg.west) -- (4.20,-3.54) |- (bmkg.east);
  \draw[wire] ([yshift=3.6mm]agg.west) -- (4.48,-3.84) |- (pvmbg.east);
  \draw[wire] ([yshift=0.6mm]agg.west) -- (4.76,-4.14) |- (pantai.east);
  \draw[wire] ([yshift=-3.1mm]agg.west) -- (4.48,-4.51) |- (feed.east);

  \draw[event] (agg) -- node[label,right]{NATS / publish} (nats);
  \draw[event] (nats.south) -- (7.15,-7.9) --
    node[label,above]{NATS / pull, ACK} (2.0,-7.9) -- (dash.north);
  \draw[event] (nats) -- node[label,right]{NATS / pull, ACK} (notify);
  \draw[event] (notify.east) -- node[label]{HTTP} (receiver.west);
\end{tikzpicture}%
}
\caption{Batas container aplikasi dan protokol komunikasi. Kotak putus-putus adalah pihak eksternal. Kotak kecil di dalam container menunjukkan penyimpanan yang dipakai oleh komponen tersebut. Container PostgreSQL memakai disk lokal dan menampung data kanonik, jurnal, dan outbox milik Aggregator. Panah berwarna mengikuti distribusi event.}
\label{fig:m2-containers}
\end{figure}


\subsection{Observabilitas}
Alur metrik, log, dan trace pada Gambar~\ref{fig:m2-telemetry} memakai container dan volume tersendiri. Telemetri dapat ditelusuri di Grafana ketika pengujian berlangsung.

\begin{figure}[H]
\centering
\resizebox{\linewidth}{!}{%
\begin{tikzpicture}[
  component/.style={draw=archink,rounded corners=2pt,fill=archlight,align=center,
    text width=3.65cm,minimum height=13mm,font=\fontsize{10}{12}\selectfont,inner sep=4pt},
  wire/.style={-{Straight Barb[length=2mm,width=1.6mm]},draw=archink,semithick},
  label/.style={fill=white,inner sep=1pt,font=\fontsize{8}{10}\selectfont,align=center}
]
  \node[component,dashed,fill=white,text width=13.1cm] (pods) at (0,0)
    {Pod aplikasi dan exporter metrik pada pod NATS};

  \node[component] (prom) at (-4.5,-2.0)
    {\textbf{Prometheus}\\Metrik pada PVC sendiri};
  \node[component] (alloy) at (0,-2.0)
    {\textbf{Grafana Alloy}\\Posisi baca pada PVC sendiri};
  \node[component] (tempo) at (4.5,-2.0)
    {\textbf{Tempo}\\Trace pada PVC sendiri};
  \node[component] (loki) at (0,-4.0)
    {\textbf{Loki}\\Log pada PVC sendiri};
  \node[component] (grafana) at (0,-6.0)
    {\textbf{Grafana}\\Konfigurasi pada PVC sendiri};

  \draw[wire] (prom.north) -- node[label,right]{HTTP /metrics} (prom.north |- pods.south);
  \draw[wire] (alloy.north) -- node[label,right]{HTTPS / log pod} (alloy.north |- pods.south);
  \draw[wire] (tempo.north |- pods.south) -- node[label,right]{OTLP/HTTP} (tempo.north);
  \draw[wire] (alloy) -- node[label,right]{HTTP / push log} (loki);
  \draw[wire] (grafana) -- node[label,right]{HTTP} (loki);
  \draw[wire] (grafana.west) -| node[label,near start,above]{HTTP} (prom.south);
  \draw[wire] (grafana.east) -| node[label,near start,above]{HTTP} (tempo.south);
\end{tikzpicture}%
}
\caption{Container observabilitas dan alur telemetri. Alloy membaca log pod melalui Kubernetes API, sedangkan aplikasi yang diinstrumentasi mengirim trace ke Tempo. Grafana mengakses tiap sumber data melalui HTTP.}
\label{fig:m2-telemetry}
\end{figure}


\section{Diagram Deployment dan Topologi}
Lingkungan acuan pengujian menggunakan Hetzner Cloud di Helsinki. OpenTofu menyiapkan empat VM Ubuntu untuk k3s dan satu VM terpisah untuk k6. \texttt{aat-dev} berperan sebagai control plane sekaligus worker. Setiap VM memiliki kernel sendiri. Penempatan VM pada host fisik dikelola penyedia. Tabel~\ref{tab:topology} mencantumkan alokasi tiap VM \cite{hetzner-cpx,hetzner-ccx}.

\begin{table}[htbp]
\centering\small
\begin{tabularx}{\textwidth}{@{}lYlrr@{}}
\toprule
VM & Peran & Tipe & vCPU & RAM \\
\midrule
\texttt{aat-dev} & Control plane dan worker & CCX33 & 8 & 32 GB \\
\texttt{aat-worker-1} & Worker & CPX32 & 4 & 8 GB \\
\texttt{aat-worker-2} & Worker & CCX23 & 4 & 16 GB \\
\texttt{aat-worker-3} & Worker & CPX32 & 4 & 8 GB \\
\texttt{aat-loadgen} & Generator beban k6 & CPX52 & 12 & 24 GB \\
\bottomrule
\end{tabularx}
\caption{Alokasi VM untuk pengujian di Hetzner Cloud.}
\label{tab:topology}
\end{table}

\begin{figure}[H]
\centering
\resizebox{\linewidth}{!}{%
\begin{tikzpicture}[
  vm/.style={draw=archink,rounded corners=3pt,text width=6.45cm,align=left,
     font=\fontsize{10}{12}\selectfont,inner sep=7pt,fill=archlight,minimum height=2.3cm},
  bus/.style={draw=archevent,rounded corners=3pt,fill=white,align=center,
     text width=8.4cm,font=\fontsize{10}{12}\selectfont,inner sep=5pt},
  link/.style={{Latex[length=1.8mm]}-{Latex[length=1.8mm]},semithick,draw=archevent}
]
  \node[vm] (c) at (-3.7,3.65) {\textbf{aat-dev}\\Control plane dan worker\\
    Client API ($\times 3$) dan Auth\\
    BMKG dan PVMBG mock\\
    Dashboard dan Notifier\\
    Penerima Tim Lapangan\\
    Grafana, operator PostgreSQL, registry\\
    Salinan shard 1, 2, 3};
  \node[vm] (w1) at (3.7,3.65) {\textbf{aat-worker-1}\\
    Aggregator dan JetStream\\
    Salinan shard 0, 2, 3};
  \node[vm,minimum height=0pt] (k6) at (3.7,6.05) {\textbf{aat-loadgen}\\Generator beban k6};
  \node[bus] (links) at (0,0) {\textbf{Jaringan privat k3s}\\
    HTTP antara API, Aggregator, dan sumber\\
    SQL/TLS ke primary dan replikasi WAL PostgreSQL\\
    Publikasi, konsumsi, dan replikasi JetStream};
  \node[vm] (w2) at (-3.7,-3.0) {\textbf{aat-worker-2}\\
    Aggregator dan JetStream\\
    Salinan shard 0, 1, 3};
  \node[vm] (w3) at (3.7,-3.0) {\textbf{aat-worker-3}\\
    Aggregator dan JetStream\\
    Salinan shard 0, 1, 2};

  \draw[-{Latex[length=1.8mm]},semithick,draw=archink] (k6.west) -|
    node[pos=0.3,above,font=\footnotesize]{HTTP / jaringan privat} (c.north);
  \draw[link] (c.south) -- (c.south |- links.north);
  \draw[link] (w1.south) -- (w1.south |- links.north);
  \draw[link] (w2.north) -- (w2.north |- links.south);
  \draw[link] (w3.north) -- (w3.north |- links.south);
\end{tikzpicture}%
}
\caption{Deployment uji beban di Hetzner. VM k6 mengakses API dan Auth dari luar klaster. Empat VM k3s menampung layanan aplikasi dan salinan shard.}
\label{fig:m2-deployment}
\end{figure}


Pada fase utama uji beban, tiga replika Client API berjalan di \texttt{aat-dev} bersama Auth, mock BMKG dan PVMBG, consumer, penerima peringatan, serta layanan observabilitas. Tiga Aggregator awal dan tiga broker JetStream tersebar di tiga worker. \texttt{aat-loadgen} mengirim permintaan kepada Client API melalui jaringan privat dan memperoleh token dari Auth.

Setiap kelompok PostgreSQL menyimpan salinannya di disk lokal pada tiga node berbeda. Label penempatan memilih node yang berhak menyimpan shard dan aturan \textit{pod anti-affinity} menjaga agar dua salinan shard yang sama tidak berada pada satu node. Ketika primary berpindah, aplikasi tetap mengakses alamat layanan untuk shard tersebut.

Seluruh pod aplikasi berada di namespace \texttt{sistem-informasi-kebencanaan-bnpb} dan diberi label \texttt{role} sesuai kelasnya. Manifest mendefinisikan \textit{resource request} dan \textit{limit}. NetworkPolicy membatasi akses langsung ke penyimpanan. Kredensial dimuat dari Kubernetes Secret dan pengaturan aplikasi dari ConfigMap.
